\section{Zmienne pomocnicze}
$a_1$ - liczba zakupionego surowca S1 powyżej 2387 ton [t] \\
$a_2$ - liczba zakupionego surowca S1 powyżej 6659 ton [t] \\
$b_1$ - liczba zakupionego surowca S2 powyżej 2090 ton [t] \\
$b_2$ - liczba zakupionego surowca S2 powyżej 4349 ton [t] \\
$c_{P1}$ - liczba wagonów transportujących surowiec S1 do przygotowalni \\ 
$e_{P2}$ - liczba ciężarówek transportujących surowiec S2 do przygotowalni \\
$e_{O2}$ - liczba ciężarówek transportującyhc surowiec S2 do zakładu obróbki cieplnej \\
$g_1$ - zmienna binarna przyjmująca wartość 1 gdy przepływ ze stanu "zakupiony surowiec S2" (S2) do stanu "surowiec S2 w zakładzie obróbki cieplnej" (O2) jest większy niż 2000 ton i mniejszy niż 4000 ton \\ 
$g_2$ - zmienna binarna przyjmująca wartość 1 gdy przepływ ze stanu "zakupiony surowiec S2" (S2) do stanu "surowiec S2 w zakładzie obróbki cieplnej" (O2) jest większy niż 4000 ton i mniejszy niż 6000 ton \\
l - liczba pracowników zatrudnionych w przygotowalni
